\section{Related Work}
\subsection{Neural Networks Pruning and Quantization}
Neural networks tend to be over-parameterized, and many algorithms have been proposed to prune away the redundant parameters. Fine-grained pruning~\cite{han2016deep, wang2020apq} cut offs the connections within the weight matrix and achieves high pruning ratios. However, it is not friendly to CPUs and GPUs and requires dedicated hardware\cite{pal2018outerspace, sparch} to support sparse matrix multiplication, which may consume extra efforts for design and automation~\cite{wang2020gcn, mao2019park, wang2018learning}. To this end, structured pruning such as channel-pruning~\cite{he2017channel, he2018amc} was further proposed to remove the entire channel to enable acceleration on general-purpose hardware.~\cite{park2016faster} further proposed to enable fine-grained pruning speedup on general-purpose platforms. However, it requires a complicated guided sparsity learning process. Quantization reduces data precision to shrink the model size and bandwidth requirements. 
\name is fundamentally different from the existing weight pruning and quantization because there is no weight in attention, and we prune the input tokens/heads based on their importance to the sentence. Quantization in \name is applied to input QKV instead of weights.


\subsection{Accelerators for Sparse and Quantized Neural Networks}
There have been various domain-specific FPGA and ASIC accelerators proposed for neural networks~\cite{mubarik2020printed, ramanathan2020look, song2020vr, kim2020duplo, liu2020duet, mo2020tfe, weng2020dsagen, vilim2020gorgon, ke2020recnmp, evans2020jpeg, gupta2020deeprecsys, zhao2020smartexchange, baek2020multi, ji2019fpsa}. Many of them leverage the sparsity to improve performance~\cite{10.1109/ISCA.2016.30, zhang2016cambricon, parashar2017scnn, zhou2018cambricon, yang2019sparse, gondimalla2019sparten, sharify2019laconic, cong2018understanding, yang2020procrustes, gong2020save, hwang2020centaur, qin2020sigma, delmas2019bit}. 
There also exist general sparse tensor algebra accelerators~\cite{chen2020tpspmv, hegde2019extensor, kanellopoulos2019smash, zhu2019sparse, kwon2019tensordimm, mahmoud2020tensordash, srivastava2020matraptor} proposed in recent years, which can be used to process sparse FC layers. 
Most of the prior work focuses on leveraging weight sparsity. By contrast, \name leverages activation (token/head) sparsity and employs specialized \topk engines to support on-the-fly cascade token/head pruning.
Quantized neural networks are also supported by many accelerators~\cite{lee2018unpu, judd2016stripes, sharify2018loom, umuroglu2018bismo, rzayev2017deeprecon, sharma2018bit, nvtensor, jouppi2017datacenter, song2020drq, zadeh2020gobo}. 
In those accelerators, the \bitwidth is fixed at the compile time, while in \name, we can adjust the \bitwidth according to the attention probability distributions.


\subsection{Efficient Natural Language Processing}
The large computation of attention-based NLP models raised much interest in improving their efficiency~\cite{hanruiwang2020hat, yan2020micronet, gu2018nonautoregressive, wang2020efficient}. GOBO~\cite{zadeh2020gobo} proposed to compress \bert model down to 3 bits, thus significantly reducing DRAM access. \powerbert~\cite{goyal2020power} prunes tokens based on the \emph{instant} attention probabilities of only one layer, which is different from \name's \emph{cumulative} attention probabilities of multiple layers. It cannot support \emph{per-head granularity} token pruning or local V vector pruning either. Head pruning is proposed in \cite{michel2019sixteen, voita2019analyzing} but they only prune head weights instead of activations as in \name. The head pruning in~\cite{voita2019analyzing} is not cascaded since pruned heads in one layer appear in latter layers. Our token pruning idea can also be generalized to Memory-Augmented Networks~\cite{NIPS2015_5846} to remove unimportant memory vectors and improve efficiency. 
